\section{Supplementary data}
\label{app:data}

The supplementary tables contain \PractitionerAccountCount{} public
first-person practitioner accounts supported by \PractitionerSourceCount{}
practitioner-source records.  Known repeated practitioners form
\PractitionerOverlapClusterCount{} overlap clusters; one controlled human--AI
workflow study is kept separate.  Accounts are project/workflow episodes rather
than unique-person records, and multiple sources may support one account.  The
collection was assembled through LLM-assisted web search and is not a systematic
or representative sample.

The screening table records \PractitionerScreenedCandidateCount{} candidates and
their dispositions.  It begins on 8 September 2026 and does not reconstruct
every source encountered in earlier searches.

The supplementary material also contains the 29-result register used for the
Davis--Kahan formalization.  It counts the four unnumbered headline theorems in
Davis--Kahan Section~2 and every named theorem, proposition, lemma, and corollary
established in the source.  Definitions, proof steps, examples, immediate
consequences, and open or deferred claims are not counted separately.
Davis--Kahan Proposition~4.4 remains a counted result because its formal
treatment gives a counterexample to the printed claim and proves a repaired
theorem.

\ifshowartifacturl
Project materials: \href{\artifacturl}{public GitHub repository}.
\else
Identifying repository URLs are suppressed in this double-blind build and are
included in the supplementary artifact.
\fi

\section{Formalization procedure}
\label{app:system}

\Cref{sec:system} summarizes the workflow.  The project used a pinned
Lean/Mathlib environment together with reusable local foundations.  Source
transcriptions, prose summaries, result records, and review judgments were kept
outside individual model conversations so work could continue across sessions.
The reviewed snapshot specifies Lean~4.34.0-rc1, Mathlib revision
\texttt{72f9c607bb3a}, and Tau Ceti revision \texttt{1b39d420ac84}.
The reusable mathematics package is \texttt{ForTauCeti}; source-facing
Davis--Kahan declarations live in the separate \texttt{DavisKahan} library.
Full dependency revisions are recorded in the supplied Lake manifest.
These identify the reproduction environment; the proposed replacement below
has not yet been validated in it.
As the project grew, we also developed tools for inspecting elaborated
declarations and checking that recorded formal evidence still referred to the
current development. Reusable foundations lived below the paper-facing theorem
modules. Agent tasks named an intended statement and its source context; agents
returned Lean declarations and obstruction notes. Compilation, axiom inspection,
and signature extraction supplied mechanical evidence, while correspondence
judgments were recorded separately. Source review included inherited section
variables and the fields of structures occurring in a theorem type, because
both can restrict a claim without appearing in its concluding inequality.

Source review compared each tracked Davis--Kahan result with the elaborated Lean
statement and definitions that affected its meaning.  Recurring questions were
the scalar field, dimension, bounded or unbounded operator scope, norm class,
gap conditions, constants, directionality, and the identity of angle and
residual operators.  Mechanical checks established properties of the formal
artifact; they did not determine whether our reading of the source was correct.

ChatGPT and Claude were used for implementation, mathematical discussion,
search, and review.  Human supervision consisted mainly of monitoring progress,
supplying framing, asking models to explain inconsistencies, and requesting
additional or independent reviews.  The humans directing the formalization were
not Davis--Kahan subject-matter experts, so disagreements about source meaning
were often investigated by probing the models in different ways rather than by
an independent expert reading of the theorem.

\section{Davis--Kahan background and Formalization 1}
\label{app:formalization1}

In the source notation, $T=A+H$ is the perturbed self-adjoint operator,
$P\mathcal H$ reduces $A$, and $Q\mathcal H$ reduces $T$
\citep{davis1970rotation}. With $E_0$ an isometry onto $P\mathcal H$,
$AE_0=E_0A_0$ defines the unperturbed trial restriction. The general residual is
\begin{equation}
 R=TE_0-E_0A_0=HE_0.
 \label{eq:dk-general-residual}
\end{equation}
In finite dimensions, the Rayleigh--Ritz choice is
\begin{equation}
 A_0=E_0^*TE_0,
 \label{eq:dk-trial-compression}
\end{equation}
which makes the trial-side diagonal block $H_0$ of the perturbation zero.
Only under this specialization is the residual exactly the component outside
the trial space:
\begin{equation}
 R=(I-P)TE_0.
 \label{eq:dk-trial-residual}
\end{equation}
It then vanishes precisely when the trial space is invariant under $T$.
In general, $R=0$ means that the specified trial operator and embedding satisfy
the exact intertwining relation, which is stronger than invariance alone.

Formalization~1 displays an \emph{ambient} theorem, not a residual theorem.
Its local names \texttt{A} and \texttt{B} denote the two bounded self-adjoint
operators; \texttt{U} reduces \texttt{A} and \texttt{V} reduces \texttt{B}.
The perturbation in its conclusion is \texttt{B - A}. No trial operator
\texttt{M} or residual \texttt{R} occurs in that declaration.

The hypotheses \texttt{hUspec} and \texttt{hUspec'} separate the restrictions
of \texttt{A} to $U$ and $U^\perp$ by a gap $d>0$: their spectra lie in
$[a,b]$ and outside $(a-d,b+d)$, respectively. This gap is on the first
operator. For the ambient double-angle bound, Davis--Kahan explicitly permit
interchanging the two operators. Thus this placement is not a defect in the
displayed ambient theorem. The directed residual inequality does not admit
that interchange in general. The historical discrepancy is the bounded
ambient operator type, not the choice of which operator supplies this gap.

For two subspaces of the same finite dimension with orthonormal coordinate
maps $E_0$ and $F_0$, their principal angles $\theta_j\in[0,\pi/2]$ satisfy
$\cos\theta_j=\sigma_j(F_0^*E_0)$.  Equivalently, they pair directions in the two
subspaces by their overlap.  The directed sine measures the part of a trial
direction outside the exact subspace; its matrix representative is
$(I-Q)E_0$.  The doubled-angle quantity applies $\sin(2\theta_j)$ to the
angles, with the trial-to-exact orientation fixed by the source.  The bound
\[
  \delta\,N(\sin 2\Theta_0)\le 2N(R)
\]
therefore controls the doubled-angle quantity by the residual relative to the
gap.  A small doubled sine alone does not imply a small angle:
$\sin(2\theta)$ vanishes at both $0$ and $\pi/2$.  Recovering a small-angle bound
requires an appropriate angle branch, such as $0\le\theta\le\pi/4$.

Formalization~1 states the estimate for complex Hilbert spaces.  Its ambient
Hilbert-space assumptions are omitted from the displayed code.  In the syntax
shown there, \texttt{E }$\mathrel{\to}\!\mathtt{L}[\mathbb C]$\texttt{ E}
means a bounded continuous complex-linear map $E\to E$.  It elaborates to
\texttt{ContinuousLinearMap (RingHom.id Complex) E E}; its boundedness is the historical scope restriction. The complex field
is not itself the project-level mismatch: a real sibling was present.  The following
table uses the same surface syntax as the displayed formalization.

\begin{table}[H]
\centering
\small
\caption{Reader's guide to the displayed ambient theorem in Formalization~1.}
\label{tab:formalization1-glossary}
\begin{tabularx}{\linewidth}{@{}>{\raggedright\arraybackslash}p{0.39\linewidth}>{\raggedright\arraybackslash}X@{}}
\toprule
Lean expression & Mathematical role \\
\midrule
\texttt{A B : E }$\mathrel{\to}\!\mathtt{L}[\mathbb C]$\texttt{ E}
  & Both ambient operators are bounded continuous complex-linear maps. \\
\texttt{IsSelfAdjoint A}
  & $A=A^*$; likewise for $B$. \\
\texttt{Reduces A U}
  & Both $U$ and $U^\perp$ are invariant under $A$; the other hypothesis
    says that $V$ reduces $B$. \\
\texttt{U.HasOrthogonalProjection}
  & The orthogonal projection onto $U$ is available; likewise for $V$.
    These instances are supplied by the surrounding context. \\
\texttt{compressOperator U A}
  & The compression of $A$ to $U$, equal to its restriction because $U$
    reduces $A$. \\
\texttt{spectrum }$\mathbb R$\texttt{ (...)}
  & The real spectrum of that self-adjoint compression. \\
\texttt{SymmetricNormingFunction}
  & The presentation's audited rename of the historical norm record.
    It supplies the ideal and gauge; the exact sidecar retains the original name. \\
\texttt{N.Mem (B - A)}
  & The perturbation lies in the ideal on which $N$ is finite. Unlike the
    current where-defined boundary, this historical theorem then \emph{concludes}
    angle ideal membership as well. \\
\texttt{N.gauge (B - A)}
  & The real-valued perturbation norm used on the right-hand side. \\
\texttt{paperSinTwoAngleOperatorC U V}
  & The ambient doubled-angle operator. Its two angle blocks must not be
    replaced by a single directed block when comparing general ideal norms. \\
\bottomrule
\end{tabularx}
\end{table}

The historical ambient conclusion asserts ideal membership and bounds the
ambient doubled-angle gauge by twice the gauge of \texttt{B - A}.  These norm and spectral conditions
are part of the statement.  The displayed declaration is complex and bounded;
a real ambient sibling also existed.  The project-level mismatch was therefore
the bounded ambient operator type, which did not cover Davis--Kahan's unbounded
setting.

\section{Common-domain double-angle replacement}
\label{app:domain-repair}

Let $A$ and $T$ be self-adjoint operators on the same dense domain
$\mathcal D\subseteq\mathcal H$. Let $P$ reduce $A$, let $Q$ reduce $T$,
and place the source gap on $T|_{Q\mathcal H}$ and $T|_{Q^\perp\mathcal H}$.
The directed clause takes a bounded map $R:P\mathcal H\to\mathcal H$ such that
\[
 Tp=Ap+Rp \qquad(p\in P\mathcal H\cap\mathcal D).
\]
The trial restriction is $A_0=A|_{P\mathcal H}$, with its own dense domain.
Since $A$ and $T$ have the same domain, $T\iota_P$ and $\iota_PA_0$ have
exactly that common domain. Neither the trial restriction nor $T-A$ is required
to be bounded. This gives the source's directed residual setup.
The ambient clause separately takes a bounded self-adjoint $H$ with $T=A+H$.
A residual hypothesis is unnecessary for that clause.

The new proof code reduces the domain issue to a bounded off-diagonal block.
Write $P^\perp=I-P$, let $\iota_P:P\mathcal H\hookrightarrow\mathcal H$, and set
\[
 C=P^\perp R\iota_P^*,\qquad S=C+C^*,\qquad J=2P-I.
\]
Reduction of $A$ and the common-domain equality imply that $P$ and $J$ preserve
$\mathcal D$. For $x\in\mathcal D$ the residual identity gives
$Cx=P^\perp TPx$. Self-adjointness of $T$, tested on the dense domain, gives
$C^*x=PTP^\perp x$. Thus
\[
 (C-C^*)x=TPx-PTx,\qquad (T-2S)Jx=JTx.
\]
The existing reflection estimate then bounds the doubled-angle block in every
Ky Fan norm by $2\|C\|_k\le 2\|R\|_k$. The existing angle correspondence and
the norm record's where-defined Fan law give the directed inequality in
\cref{eq:sin2theta}. The ambient clause uses the separate existing ambient
theorem. The proposed Lean module follows these steps; its compiler validation
remains pending. These equations explain the proposed repair, not a claim that
the new Lean declarations have already been checked.

The restriction in the current theorem is also visible with a finite gap
interval. Take $A=T=D\oplus0$ on $\ell^2\oplus\ell^2$, with
$D(x_n)=((n+1)x_n)$, let $P$ select the first summand and $Q$ the second.
Then $\sigma(T|_Q)=\{0\}$, $\sigma(T|_{Q^\perp})\subseteq[1,\infty)$,
and the gap condition holds with $[\beta,\alpha]=[0,0]$, $\delta=1$.
The source's dimension conditions hold and $H=R=0$. The angle is $\pi/2$,
so its doubled sine is zero; $P\mathcal H$ still fails the whole-space domain
hypothesis. This example also shows why a small doubled sine alone does not
select the small-angle branch.

The residual-only boundedness issue is independent. With
$A=(-D)\oplus D$, $T=(-D)\oplus2D$, and $P=Q$ the first summand,
the operators have the same domain, $R=0$, and the perturbed spectral parts
are separated. The difference $T-A=0\oplus D$ is unbounded. The directed
source inequality remains applicable, although an endpoint requiring a
bounded ambient perturbation cannot be instantiated.

\paragraph{Norm representation.}
The current norm record contains a symmetric ideal gauge, rank-one
normalization, and the comparison implication
\[
 \bigl(\forall k,\ \|X\|_k\le\|Y\|_k\bigr)
 \ \Longrightarrow\ N(X)\le N(Y)
\]
when both $X$ and $Y$ belong to the norm's domain. The theorem exposing this
property projects the corresponding structure field. This uses the
source-cited Fan comparison as packaged mathematical background; it is not an
independent formal derivation of Fan dominance from the bare norm axioms.
The new common-domain code retains that boundary and makes no additional
norm-representation claim.

\paragraph{Transcription check.}
The supplied modernized transcription states the factor $2$ in the headline
residual inequality but omits it in the residual step of the Section~7 proof.
The factor $2$ is necessary. For
\[
 A=0,\quad T=H=\begin{pmatrix}0&1\\1&0\end{pmatrix},\quad
 P\mathbb R^2=\operatorname{span}(e_1),\quad
 Q\mathbb R^2=\operatorname{span}(e_1+e_2),
\]
the spectral gap is $2$, the angle is $\pi/4$, and $\|R\|=1$.
Hence $\delta\|\sin2\Theta_0\|=2=2\|R\|$; the inequality with factor $1$
is false. The reflection calculation also gives $H-JHJ$ equal to twice the
off-diagonal part. This identifies an inconsistency in the supplied
transcription; locating its origin requires comparison with the printed scan.

\section{Lean publication activity}

\begin{figure}[H]
  \centering
  \input{generated/lean_activity_timeline_tikz.tex}
  \caption{Monthly publication counts reproduced from the Papers With Lean
  public statistics series, January 2025 through August 2026
  \citep{paperswithlean2026stats}.  Counts are grouped by the captured data's
  \texttt{published} month, and only complete months are shown.  The fraction
  that used AI is unknown.}
  \label{fig:lean-activity}
\end{figure}



\section{Semantic-review page}
\label{app:review-page}

\IfFileExists{figures/semantic-alignment-sine-theta-row.png}{
\begin{figure}[H]
  \centering
  \includegraphics[width=\linewidth,trim=0 250 0 0,clip]{figures/semantic-alignment-sine-theta-row.png}
  \caption{A review interface developed late in the project.  The source passage
  and inherited context appear at left, the claimed source-to-Lean
  correspondence in the middle, and compiler-derived Lean evidence at right.
  Most reviews reported in the main paper used LLM source/signature comparison;
  the interface was added later.}
  \label{fig:dashboard}
\end{figure}

}{
The dashboard image is omitted from this review build pending regeneration.
It is not used as evidence for the current source-correspondence judgment.
}

The interface consolidates source passages, correspondence judgments, and
compiler-derived signatures previously inspected through files and Lean
queries.  Understanding a signature can also require the definitions of its
structures, predicates, coercions, and typeclass assumptions.  Lean Compass's
semantic dependency view suggests how these definitions could be exposed while
omitting proof-only dependencies \citep{yanahama2026leanatlas}.

\section{Resource utilization}
\label{app:resources}

To estimate the cost of the formalization effort, we maintained a resource
accounting record.  The compact snapshot below has an accounting cutoff of
13 August 2026.  Retained model turns and token quantities are incomplete
measured lower bounds; price and serving-energy rows are estimates derived from
those observations.  The energy estimate excludes training, embodied hardware,
client devices, and network energy.

\begin{table}[H]
\centering
\scriptsize
\caption{Observed lower bounds and derived estimates at the 13 August 2026
accounting cutoff.  These are not project totals.}
\begin{tabularx}{\linewidth}{@{}p{0.28\linewidth}p{0.17\linewidth}>{\raggedright\arraybackslash}X@{}}
\toprule
Quantity & Lower bound / estimate & Basis / interpretation \\
\midrule
\input{generated/resource_snapshot_table.tex}
\end{tabularx}
\end{table}


\begin{table}[H]
\centering
\scriptsize
\caption{AI systems and model families used during the formalization.
Token telemetry is incomplete.}
\label{tab:formalization-systems}
\begin{tabularx}{\linewidth}{@{}p{0.24\linewidth}>{\raggedright\arraybackslash}X@{}}
\toprule
System & Model families observed \\
\midrule
\input{generated/model_systems_table.tex}
\end{tabularx}
\end{table}

\begin{table}[H]
\centering
\scriptsize
\caption{Recovered model-resolved token measurements at the 13 August 2026
accounting cutoff.  These are measured lower bounds for the five labels with
recovered telemetry; models in \cref{tab:formalization-systems} that do not
appear here were still used.}
\label{tab:model-tokens-appendix}
\resizebox{\linewidth}{!}{% Generated by scripts/build_accounting.py; do not edit by hand.
\begin{tabular}{lrrrr}
\toprule
Model & \multicolumn{4}{c}{Retained measured tokens} \\
 & Input & Cache write & Cache read & Output \\
\midrule
claude-opus-5 & 140,093 & 116,791,540 & 17,395,271,984 & 34,338,937 \\
claude-opus-4-8 & 693,366 & 58,990,139 & 3,744,966,700 & 15,967,324 \\
claude-fable-5 & 380,759 & 27,789,153 & 1,512,917,593 & 7,985,843 \\
gpt-5.6-sol & 6,502,664 & 0 & 343,618,944 & 635,601 \\
gpt-5.6-terra & 3,798,915 & 0 & 382,951,936 & 563,722 \\
claude-sonnet-5 & 78 & 110,211 & 2,256,550 & 26,074 \\
\bottomrule
\end{tabular}
}
\end{table}

Historical model labels are retained verbatim, with exact recovered telemetry used
when it is available.
